Los histogramas constituyen la base de varias técnicas de procesamiento
en el dominio espacial. La manipulación de los histogramas es usada de
manera eficiente en el realce o mejoramiento de la calidad de una imagen.
La información estadística obtenida a partir de los histogramas se utiliza
en diversas aplicaciones como compresión y segmentación de imágenes. La
facilidad con la que se pueden calcular los histogramas usando software y
su bajo consumo de recursos de hardware en su implementación, han hecho de
esta herramienta una de las más usadas en el procesamiento en tiempo real.

El histograma de una imagen es la representación gráfica de la distribución
que existe de las distintas tonalidades de grises con relación al número de
píxeles o porcentaje de los mismos, es decir, un histograma representa la
frecuencia relativa de ocurrencia de los niveles de gris. La representación
de un histograma ideal sería la de una recta horizontal, ya que eso nos
indicaría que todos los posibles valores de grises están distribuidos de
manera uniforme en nuestra imagen.

La ecualización del histograma es una técnica bastante conocida y sirve para
obtener un histograma uniforme de tal manera que los niveles de gris son
distribuidos sobre la escala y un número igual de píxeles son colocados en
cada nivel de gris. Para un observador, esta ecualización hace que las imágenes
se vean más balanceadas y con mejor contraste. Como consecuencia, una imagen
ecualizada, permite que ciertos detalles sean visibles en regiones oscuras o
brillantes.

Un pixel $p$ en las coordenadas $(x, y)$ tiene 4 vecinos horizontales y 4
verticales cuyas coordenadas están dadas por:
\begin{equation}
    (x+1,y), (x-1,y), (x,y+1), (x,y-1)
\end{equation}

Este grupo de píxeles se anota con $N_{4}(p)$. Así mismo, las vecindades
diagonales con el punto $(x,y)$ se anotan como $N_{D}(p)$ y sus coordenadas son:
\begin{equation}
    (x+1,y+1), (x-1, y-1), (x-1,y+1), (x+1, y-1)
\end{equation}

Para definir de forma adecuada el concepto de vecindad, es necesario revisar el
de adyacencia. Dos píxeles son adyacentes si, y sólo si, tienen en común una de
sus fronteras, o al menos una de sus esquinas. La figura 1 muestra píxeles adyacentes.

Dos píxeles son vecinos si cumplen con la definición de adyacencia. Si los
píxeles comparten una de sus fronteras, se dice que los mismos son vecinos
directos; si solo se tocan en una de sus esquinas, se llaman vecinos indirectos.

Una vecindad de un pixel $p_{0}$ denota como $V_{p}$, es una submatriz $M_{KL}$
de tamaño $K \times L$ con $K$ y $L$ enteros impares pequeños, contenida en la
matriz imagen $(i_{MN})$, la cual está formada por un número finito de pixeles
vecinos o no de $p_{0}$.

\begin{equation}
    V_{p}=\{p:p \in M_{KL}\}; \quad M_{KL} \subset i_{MN}; \quad K=L=3,5,9
\end{equation}

Para ilustrar las definiciones anteriores, se puede observar en la figura 2,
vecindades de 4 y de vecindades 8, la primera formada por píxeles que son vecinos
directos, mientras que la vecindad de 8 está formada tanto por vecinos directos
como indirectos.

En el desarrollo de las técnicas de procesamiento de imágenes que involucren el
análisis de una determinada región de la escena digital, es posible encontrar
vecindades de $5 \times 5$ y hasta $9 \times 9$ básicamente la definición de las
dimensiones de la matriz vecindad depende de la técnica que se esté desarrollando.

\subsection*{Principio del algoritmo LBP}
El algortimo de patrones binarios locales (LBP por sus siglas en ingles Local
Binary Pattern) es un algoritmo muy usado en visión por computadora para describir
la textura de una imagen. Funciona comparando la intensidad de un píxel con la de
sus vecinos y generando un código binario que representa el patrón local. El
algoritmo de LBP calcula una puntuación para cada píxel comparándolo con sus vecinos.

Matemáticamente, para un píxel central $g_{c}$ y $P$ vecinos en un radio $R$:
\begin{equation}
    LBP_{P,R}=\sum_{p=0}^{P-1}s(g_{p}-g_{c})2^{p}
\end{equation}

Donde $s(x)$ es la función signo:
\begin{itemize}
    \item $s(x)=1$ si $x \ge 0$
    \item $s(x)=0$ si $x < 0$
\end{itemize}

Para entender la técnica de LBP, debemos mirar dentro de una ventana de $3 \times 3$
píxeles. El algoritmo convierte la textura local en un número entero simple.

\textbf{Mecánica de Cálculo Manual (Invariante a la Rotación):}
\begin{enumerate}
    \item \textbf{Umbralización:} Compara cada uno de los 8 vecinos $(P=8, R=1)$ con
          el píxel central. Si el vecino es mayor o igual, se asigna '1'; de lo contrario, '0'.
    \item \textbf{Ponderación y Suma:} Los bits resultantes se multiplican por
          potencias de 2 $(2^{0}, 2^{1}, ..., 2^{7})$ siguiendo un orden (por ejemplo,
          sentido horario comenzando arriba a la izquierda) y se suman para obtener el
          valor LBP decimal.
\end{enumerate}